\chapter{Ergodicity Economics}
\label{ch:ergodicity}

\epigraph{\itshape Three hundred years of expected utility theory
rested on a confusion of two very different averages.}{--- Paraphrase
of Peters \& Gell-Mann (2016)}

\noindent
This chapter develops the second major resource of Part II:
ergodicity economics. The treatment is the most directly applied of
any in this part of the book; the framework we develop here is used
in essentially every chapter of Part III, but most decisively in
Chapter 12 (halal versus haram dynamics) and Chapter 13 (riba as
financial entropy).

The chapter assumes some familiarity with intermediate economics ---
the concepts of expected utility, risk aversion, and basic finance
should be familiar. It assumes no background in measure-theoretic
probability or stochastic processes; both are developed where
needed. The mathematical level is set so that a graduate student in
economics can follow without further preparation.

The chapter has two main goals. The first is to develop the formal
machinery of ergodicity --- random variables, expectations, time
versus ensemble averages --- with the precision required for the
worked treatments in Part III. The second is to demonstrate the
specific empirical and theoretical claims that ergodicity economics
makes about wealth dynamics, and to begin to draw out the
implications for Islamic economic claims.

The chapter proceeds in eleven sections.
\xref{sec:probability-basics} establishes the basic vocabulary of
probability theory.
\xref{sec:bernoulli-puzzle} introduces the St.\ Petersburg paradox
and the Bernoulli framework that has shaped expected utility theory.
\xref{sec:hidden-assumption} identifies the hidden assumption ---
ergodicity --- in the standard framework.
\xref{sec:ensemble-time} develops the formal distinction between
ensemble and time averages.
\xref{sec:coin-flip} works the canonical multiplicative coin-flip
example in full.
\xref{sec:multiplicative-dynamics} develops multiplicative wealth
dynamics in general.
\xref{sec:kelly-criterion} treats the Kelly criterion.
\xref{sec:absorbing-barriers} develops absorbing barriers and their
implications.
\xref{sec:dissolved-puzzles} shows how several long-standing economic
puzzles are dissolved by the ergodicity framework.
\xref{sec:risk-pooling} treats risk pooling as the rational response
to non-ergodic dynamics.
\xref{sec:riba-preview} previews the connection to riba, developed
in Chapter 13.

\section{Probability foundations}
\label{sec:probability-basics}

We begin with the basic vocabulary, stated precisely so that later
sections can build on it. The reader who has studied probability
formally can skim; the reader who has not should attend.

\subsection{Random variables and probability spaces}

A \emph{probability space} is a triple $(\Omega, \mathcal{F}, P)$ where:

\begin{itemize}[leftmargin=2em]
    \item $\Omega$ is the \emph{sample space}: the set of all
    possible outcomes of the random experiment.
    \item $\mathcal{F}$ is a \emph{$\sigma$-algebra} (sigma-algebra)
    on $\Omega$: a collection of subsets of $\Omega$ closed under
    complements and countable unions, called \emph{events}.
    \item $P$ is a \emph{probability measure}: a function from
    $\mathcal{F}$ to $[0, 1]$ such that $P(\Omega) = 1$ and
    countable disjoint unions are additive.
\end{itemize}

A \emph{random variable} $X$ is a measurable function from $\Omega$
to the real numbers $\mathbb{R}$. The reader can think of $X$ as a
quantity whose value depends on the random outcome: when the
experiment is run and outcome $\omega \in \Omega$ occurs, the random
variable takes the value $X(\omega)$.

\begin{example}[A coin flip]
Let $\Omega = \{H, T\}$ be the sample space for a single coin flip.
Let $\mathcal{F} = \{\emptyset, \{H\}, \{T\}, \Omega\}$ be the
$\sigma$-algebra. Let $P(\{H\}) = P(\{T\}) = 1/2$ for a fair coin.
Define a random variable $X$ by $X(H) = 1, X(T) = -1$. Then $X$ is
a real-valued random variable taking values $+1$ or $-1$ with equal
probability.
\end{example}

\subsection{Expectation, variance, and moments}

The \emph{expected value} of a random variable $X$ is
\[
    \mathbb{E}[X] = \int_\Omega X(\omega)\, dP(\omega).
\]
For a discrete random variable taking values $x_1, x_2, \dots, x_n$
with probabilities $p_1, p_2, \dots, p_n$, this becomes
\[
    \mathbb{E}[X] = \sum_{i=1}^n p_i x_i.
\]
For a continuous random variable with density $f_X(x)$,
\[
    \mathbb{E}[X] = \int_{-\infty}^\infty x f_X(x)\, dx.
\]

The expected value is a single real number that we interpret as the
``average'' value of $X$. The interpretation is not innocent, as we
shall see; it conflates two distinct notions of ``average'' that
turn out to be different in important cases.

The \emph{variance} of $X$ is
\[
    \text{Var}(X) = \mathbb{E}[(X - \mathbb{E}[X])^2] =
    \mathbb{E}[X^2] - (\mathbb{E}[X])^2.
\]
The variance measures the spread of $X$ around its mean. Higher
moments ($\mathbb{E}[X^k]$ for $k \geq 3$) measure further features
of the distribution; we will use them sparingly.

\subsection{The law of large numbers}

The fundamental result of classical probability theory, due to Jacob
Bernoulli (1713):

\begin{theorem}[Strong Law of Large Numbers]
\label{thm:lln}
Let $X_1, X_2, X_3, \dots$ be a sequence of independent and
identically distributed random variables with finite mean
$\mathbb{E}[X_i] = \mu$. Then with probability 1,
\[
    \lim_{n \to \infty} \frac{1}{n} \sum_{i=1}^n X_i = \mu.
\]
\end{theorem}

In words: if you repeat the experiment many times and take the
sample average of the results, the sample average converges to the
expected value as the number of repetitions grows.

The Law of Large Numbers is the formal foundation for the intuitive
identification of ``expected value'' with ``long-run average.'' If
you flip a fair coin many times and count the proportion of heads,
the proportion approaches 1/2.

But notice the structure of the theorem carefully. The
``many-times'' is across many \emph{independent repetitions}. The
sample average is the average across many copies. The theorem says
that this average converges to the expected value.

What if we cannot get many independent copies? What if the
experiment we care about is something that happens once, sequentially,
to a single agent over time? Then the Law of Large Numbers does not
directly apply. The interpretation of expected value as ``long-run
average'' assumes that ``long-run'' means many independent copies, not
many sequential trials by a single agent.

For most economic applications, the relevant ``long run'' is
sequential, not parallel. An individual investor's wealth dynamics
unfold over time, in a single trajectory. The investor cannot run
many copies of their own life in parallel. The Law of Large Numbers
does not, in this case, validate the use of expected value as a
decision criterion.

This is the central observation of ergodicity economics. We unpack
it carefully in the next several sections.

\section{The Bernoulli framework and the St.\ Petersburg paradox}
\label{sec:bernoulli-puzzle}

\subsection{The St.\ Petersburg paradox}

In 1738 Daniel Bernoulli (Jacob's nephew) published an analysis of a
gambling problem proposed by his cousin Nicolas, now known as the
\emph{St.\ Petersburg paradox}.\footnote{Bernoulli, D. (1738),
``Specimen theoriae novae de mensura sortis'' (``Exposition of a new
theory on the measurement of risk''), translated by Sommer (1954),
\emph{Econometrica} 22: 23--36.}

The setup: a fair coin is flipped repeatedly until it lands heads.
If the first heads occurs on the $k$-th flip, the player wins
$2^k$ ducats. The expected payoff is
\[
    \mathbb{E}[X] = \sum_{k=1}^\infty \frac{1}{2^k} \cdot 2^k =
    \sum_{k=1}^\infty 1 = \infty.
\]
The expected value is infinite. By standard expected-value
reasoning, the player should be willing to pay any finite amount to
play this game. Yet, asked how much they would actually pay, real
people offer modest amounts --- a few ducats, perhaps a few dozen.

The discrepancy is the paradox. Either people are irrational (in
which case standard decision theory misdescribes their behavior), or
expected value is the wrong decision criterion (in which case
standard decision theory is missing something).

\subsection{Bernoulli's resolution}

Bernoulli proposed that what people maximize is not expected
\emph{wealth} but expected \emph{utility} of wealth, where utility
is a concave function of wealth. The most common choice is the
logarithm:
\[
    \mathbb{E}[\log W] = \sum_{k=1}^\infty \frac{1}{2^k} \log(W_0 + 2^k - c),
\]
where $W_0$ is the player's initial wealth and $c$ is the price of
playing. Because $\log$ grows much more slowly than linearly, the
sum converges, and a finite price equates expected utility with the
status quo.

This is the foundation of \emph{expected utility theory}, formalized
by von Neumann and Morgenstern (1944) and now the standard framework
in economics and finance.\footnote{von Neumann, J. \& Morgenstern,
O. (1944), \emph{Theory of Games and Economic Behavior}, Princeton
University Press.} Modern decision theory teaches that rational
agents maximize the expected value of a utility function over
outcomes; the curvature of the utility function captures their
attitude toward risk.

The framework has become so standard that few economists ask whether
the underlying maneuver --- using \emph{expected} utility ---
involves a substantive assumption. It does. The assumption is that
the relevant criterion is the ensemble average of the utility, which
is the same as the time average only if the process is ergodic.

\subsection{Why Bernoulli's resolution is incomplete}

Bernoulli's introduction of utility was a real advance, but it
missed the deeper issue. The expected-utility framework still
computes an \emph{expectation} (the ensemble average), then
interprets that expectation as guiding individual decisions. The
question of whether the ensemble average represents the experience
of the individual is not asked.

Ergodicity economics, beginning with the work of Peters (2011) and
Peters and Gell-Mann (2016), has shown that the question must be
asked. For non-ergodic processes --- including most realistic
wealth dynamics --- the ensemble average and the time average
diverge. The agent who maximizes the ensemble average makes
decisions appropriate for an ensemble; the agent who experiences
their own wealth over time makes decisions on the time average. The
two are different criteria with different optimal strategies.

We develop this in the rest of the chapter.

\section{The hidden assumption}
\label{sec:hidden-assumption}

We now state the hidden assumption explicitly.

\begin{conceptbox}[The hidden assumption of standard decision theory]
\label{conc:hidden-assumption}
Standard decision theory uses the expected value (ensemble average)
of the relevant random variable as the criterion for individual
decisions. This is appropriate if and only if the relevant process
is ergodic --- that is, if the time average experienced by an
individual realization equals the ensemble average across
realizations.

For non-ergodic processes, the ensemble average is the wrong
criterion for the individual. The individual experiences the time
average; their decisions, if rational, should optimize the time
average.

The assumption of ergodicity is rarely made explicit in standard
decision theory, but it is required to justify the use of expected
value as a decision criterion.
\end{conceptbox}

The reader unfamiliar with this distinction may find the claim
disorienting. The ensemble average and the time average are
\emph{usually thought to be the same thing}; the conflation is
built into intuitive thinking about probability and into most
treatments of decision theory.

The conflation is correct for ergodic processes, which include some
important cases (notably, the additive random walks and stationary
processes that the Law of Large Numbers most directly addresses).
But it is incorrect for non-ergodic processes, which include
multiplicative wealth dynamics, processes with absorbing barriers
(such as bankruptcy), and many other realistic cases.

The next section develops the distinction with mathematical
precision.

\section{Ensemble and time averages: definitions and the difference}
\label{sec:ensemble-time}

\subsection{Two kinds of average}

Suppose we have a stochastic process $X(t)$ defined for $t \geq 0$.
Two natural ways to compute an ``average'' of the process:

\begin{definition}[Ensemble average]
\label{def:ensemble-avg}
The \emph{ensemble average} of $X$ at time $t$ is
\[
    \langle X \rangle_{\text{ens}}(t) = \mathbb{E}[X(t)] =
    \int X(t, \omega)\, dP(\omega),
\]
the average of $X(t)$ across all possible realizations of the
process. If we run the process many independent times, the
ensemble average is the limit of the sample average across runs.
\end{definition}

\begin{definition}[Time average]
\label{def:time-avg}
The \emph{time average} of $X$ for a single realization is
\[
    \langle X \rangle_{\text{time}} = \lim_{T \to \infty}
    \frac{1}{T} \int_0^T X(t)\, dt,
\]
the average of $X$ over time within a single realization (provided
the limit exists).
\end{definition}

For multiplicative processes, the relevant time average is not the
arithmetic time average but the time-averaged growth rate:
\[
    g_{\text{time}} = \lim_{T \to \infty} \frac{1}{T} \log\frac{X(T)}{X(0)}.
\]
This is the quantity that determines the long-run growth of an
individual realization.

\subsection{The definition of ergodicity}

\begin{definition}[Ergodicity]
\label{def:ergodicity}
A process $X(t)$ is \emph{ergodic} (with respect to the relevant
average) if the time average for a single realization equals the
ensemble average:
\[
    \langle X \rangle_{\text{time}} = \langle X \rangle_{\text{ens}}
    \quad \text{(with probability 1).}
\]
A process for which this equality fails is \emph{non-ergodic}.
\end{definition}

The textbook examples of ergodic processes include independent
identically distributed sums (the Law of Large Numbers result) and
many stationary stochastic processes. The textbook examples of
non-ergodic processes include multiplicative wealth dynamics,
processes with absorbing barriers, and dynamical systems with
multiple equilibrium attractors.

The terminology comes from Boltzmann's work on statistical mechanics
in the late nineteenth century. Boltzmann assumed that physical
systems are ergodic --- that the time average of a quantity over a
long enough time equals the ensemble average over a phase-space
distribution --- in order to derive the laws of thermodynamics from
microscopic mechanics. The assumption is satisfied for many physical
systems but not all. For our purposes, what matters is that economic
systems involving wealth dynamics are typically \emph{not} ergodic.

\subsection{Why ensemble and time can differ}

A simple example illustrates how the two averages can diverge.
Consider a two-step random process. At each step, the value $X$ is
multiplied by either 2 (probability 1/2) or 1/2 (probability 1/2).
Starting from $X(0) = 1$.

\textbf{Ensemble computation.} After one step, the expected value of
$X$ is
\[
    \mathbb{E}[X(1)] = \frac{1}{2} \cdot 2 + \frac{1}{2} \cdot \frac{1}{2}
    = 1.25.
\]
After $n$ steps, by independence,
\[
    \mathbb{E}[X(n)] = (1.25)^n.
\]
The ensemble average grows exponentially, with growth rate
$\log(1.25) \approx 0.223$ per step.

\textbf{Time computation.} For a single realization, after $n$
steps with $k$ multiplications by 2 and $n-k$ multiplications by 1/2,
\[
    X(n) = 2^k \cdot (1/2)^{n-k} = 2^{2k-n}.
\]
With probability 1/2 of multiplying by 2 each step, $k \approx n/2$
for large $n$ (by the Law of Large Numbers applied to the
\emph{number of} multiplications). So for a typical realization,
\[
    X(n) \approx 2^{n - n} = 1
\]
or more precisely, the time-averaged growth rate is
\[
    g_{\text{time}} = \lim_{n \to \infty} \frac{1}{n} \log X(n)
    = \frac{1}{2} \log 2 + \frac{1}{2} \log\frac{1}{2} = 0.
\]
A typical realization has zero long-run growth rate.

\textbf{The discrepancy.} The ensemble average grows at rate $\log
1.25$. The time average grows at rate 0. These are different. Both
are correct; they are answering different questions. The ensemble
average answers ``what is the average behavior across many
independent realizations?''; the time average answers ``what is the
behavior of a typical individual realization?''

For an investor facing this gamble, the relevant question is the
second. The investor experiences a single realization, not the
ensemble. They will, with high probability, end up with $X(n)$ near
1, not near $(1.25)^n$. The ensemble average is, in this sense,
misleading about individual experience.

\section{The canonical multiplicative coin-flip example}
\label{sec:coin-flip}

We now develop the canonical example of ergodicity economics in
full. The example is structurally simple but mathematically
illuminating.

\subsection{Setup}

A fair coin is flipped repeatedly. Each flip:

\begin{itemize}[leftmargin=2em]
    \item If heads (probability 1/2), wealth is multiplied by
    $1 + a$ for some positive $a$.
    \item If tails (probability 1/2), wealth is multiplied by
    $1 - b$ for some positive $b$, with $b < 1$.
\end{itemize}

Starting wealth is $W_0$. After $n$ flips with $h$ heads and $t = n - h$
tails:
\[
    W(n) = W_0 (1+a)^h (1-b)^{n-h}.
\]

\subsection{Ensemble average}

The expected wealth after one flip is
\[
    \mathbb{E}[W(1)] = W_0 \left[ \frac{1}{2}(1+a) + \frac{1}{2}(1-b) \right]
    = W_0 \left[ 1 + \frac{a-b}{2} \right].
\]
Define $r_e = \frac{a-b}{2}$, the \emph{ensemble-average return per
flip}. Then
\[
    \mathbb{E}[W(n)] = W_0 (1 + r_e)^n.
\]
If $a > b$, the ensemble average grows; if $a < b$, it shrinks; if
$a = b$, it stays the same.

\subsection{Time average}

For large $n$, by the Law of Large Numbers, $h \approx n/2$ and
$n - h \approx n/2$. The realized wealth is approximately
\[
    W(n) \approx W_0 (1+a)^{n/2} (1-b)^{n/2} = W_0 \left[ (1+a)(1-b) \right]^{n/2}.
\]
The time-averaged growth rate is
\[
    g_{\text{time}} = \lim_{n \to \infty} \frac{1}{n} \log\frac{W(n)}{W_0}
    = \frac{1}{2} \log[(1+a)(1-b)].
\]
Define $r_t$ implicitly by $1 + r_t = e^{g_{\text{time}}} =
\sqrt{(1+a)(1-b)}$. This is the \emph{time-average return per flip}.

\subsection{Comparison: the +50\%/-40\% case}

Consider the specific case $a = 0.5, b = 0.4$. So heads multiplies
wealth by 1.5, tails by 0.6.

\textbf{Ensemble average.}
\[
    r_e = \frac{0.5 - 0.4}{2} = 0.05.
\]
The ensemble-average return is $+5\%$ per flip. Standard decision
theory, applied naively, would say: this is a profitable gamble, the
investor should play it.

\textbf{Time average.}
\[
    g_{\text{time}} = \frac{1}{2} \log[(1.5)(0.6)] = \frac{1}{2} \log(0.9) \approx -0.0527.
\]
The time-average growth rate is approximately $-5.3\%$ per flip.
The investor's wealth shrinks over time at this rate.

\textbf{The startling discrepancy.} In the ensemble, the gamble has
positive expected return: across many investors, the average wealth
grows. In the individual time series, the gamble has negative
return: the typical investor's wealth shrinks.

The two statements are both correct. They describe different
quantities. An investor who follows the ensemble-average advice and
plays this gamble will, with probability approaching 1, lose
substantial wealth over time. The ``positive expected return'' is
real but it lives in a few extreme realizations where one happens
to flip heads many more times than tails; in the typical
realization, the wealth shrinks.

\subsection{Why this happens}

The mechanism is geometric. Multiplicative random factors satisfy
the inequality
\[
    \text{geometric mean} \leq \text{arithmetic mean}
\]
with equality only when all factors are equal. The arithmetic mean
of the two factors $\{1.5, 0.6\}$ is $1.05$ (the ensemble average).
The geometric mean is $\sqrt{1.5 \cdot 0.6} = \sqrt{0.9} \approx 0.949$
(below 1). The arithmetic mean is what the ensemble grows at; the
geometric mean is what the typical individual experiences.

This is a mathematical fact about multiplicative dynamics. It has
nothing to do with risk aversion, behavioral biases, or
psychological idiosyncrasy. It is structural.

\subsection{The general principle}

Generalize: for any random multiplicative factor $1 + R$ where $R$
is a non-degenerate random variable with $\mathbb{E}[R] > 0$ and
$\text{Var}(R) > 0$, the long-run time-averaged growth rate
satisfies
\[
    g_{\text{time}} = \mathbb{E}[\log(1+R)] < \log(1 + \mathbb{E}[R])
    = g_{\text{ens}}
\]
by Jensen's inequality (since $\log$ is concave). The strict
inequality means the time average is always strictly less than the
ensemble average for any non-trivial random multiplicative factor.

\begin{theorem}[Time-ensemble divergence for multiplicative dynamics]
\label{thm:time-ensemble-divergence}
Let $W_0, W_1, W_2, \dots$ be a wealth process with
$W_{n+1} = W_n (1 + R_{n+1})$ where the $R_n$ are independent and
identically distributed random variables with finite second moment.
Let $r = \mathbb{E}[R_n]$. The ensemble-average growth rate is
$\log(1 + r)$. The time-averaged growth rate is
$\mathbb{E}[\log(1 + R_n)]$. By Jensen's inequality,
\[
    \mathbb{E}[\log(1 + R_n)] \leq \log(1 + r),
\]
with equality only if $R_n$ is degenerate (constant). Strict
inequality holds whenever $R_n$ has positive variance.
\end{theorem}

\begin{proof}
The function $\log(1+x)$ is strictly concave for $x > -1$. By
Jensen's inequality, for any non-degenerate random variable $X$
taking values in $(-1, \infty)$,
\[
    \mathbb{E}[\log(1+X)] < \log(1 + \mathbb{E}[X]).
\]
Apply this to $X = R_n$. The conclusion follows.
\end{proof}

The theorem is the formal foundation of all the qualitative claims
of ergodicity economics. The ensemble average always overstates the
time-averaged growth for multiplicative dynamics with random
returns. The standard expected-value framework, applied without
recognizing this, systematically advises strategies that reduce
individual long-run wealth.

\section{The Kelly criterion}
\label{sec:kelly-criterion}

Given that the time-averaged growth rate is the appropriate
criterion for an investor experiencing wealth over time, what
investment fraction maximizes it? The answer was given by John Kelly
in 1956, in a paper that has subsequently become foundational for
modern probability theory and information theory.\footnote{Kelly,
J. L. (1956), ``A new interpretation of information rate,''
\emph{Bell System Technical Journal} 35: 917--926.}

\subsection{The setup}

Consider a sequence of bets, each with probability $p$ of winning and
$q = 1 - p$ of losing. A win returns $b$ times the amount bet (so
betting \$1 returns \$$b$ on a win, plus the original \$1; a loss
loses the entire bet).

The investor begins with wealth $W_0$ and bets a fraction $f$ of
their current wealth on each round. After a win, wealth becomes
$W(1 + bf)$. After a loss, wealth becomes $W(1 - f)$.

The objective: choose $f$ to maximize the long-run time-averaged
growth rate of wealth.

\subsection{The Kelly formula}

The time-averaged growth rate as a function of $f$ is
\[
    g(f) = p \log(1 + bf) + q \log(1 - f).
\]
To maximize, take the derivative with respect to $f$ and set it
equal to zero:
\[
    \frac{dg}{df} = \frac{pb}{1 + bf} - \frac{q}{1 - f} = 0.
\]
Solving:
\[
    pb(1 - f) = q(1 + bf),
\]
\[
    pb - pbf = q + qbf,
\]
\[
    pb - q = (pb + q) bf,
\]
\[
    f^* = \frac{pb - q}{b(pb + q)} = \frac{p - q/b}{1 + q b / b} = p - \frac{q}{b}.
\]
The Kelly fraction is
\[
    \boxed{f^* = p - \frac{q}{b}.}
\]

\subsection{What the formula says}

The Kelly fraction depends on both the win probability $p$ and the
payoff ratio $b$. Several special cases:

\begin{itemize}[leftmargin=2em]
    \item If $pb \leq q$ (the expected gross return is non-positive),
    then $f^* \leq 0$: don't bet (or bet against, if shorting is
    available).
    \item If $b = 1$ (even-money bet) and $p = 0.6, q = 0.4$, then
    $f^* = 0.6 - 0.4 = 0.2$: bet 20\% of wealth.
    \item If $b$ is very large (high payoff), $q/b \to 0$ and
    $f^* \to p$: bet the win probability.
    \item If $p$ is close to 1 (near-certain win), $f^*$ is close
    to 1 minus a small correction: bet most of wealth.
\end{itemize}

\subsection{What the formula does not say}

The Kelly criterion is not the same as expected-value maximization.
A naive expected-value maximizer who computed
\[
    \mathbb{E}[\Delta W \mid \text{bet fraction } f]
    = pf b - qf > 0 \quad \text{whenever } pb > q
\]
would conclude: bet everything. The maximum expected return
strategy is $f = 1$. But $f = 1$ guarantees ruin on the first loss.
The Kelly fraction is strictly less than 1 except for the
degenerate case of certain wins; it represents the intermediate
between cowardly under-betting and ruinous over-betting that
maximizes long-run growth.

The Kelly criterion is an example of \emph{time-average optimization}
applied correctly. It tells the agent: to maximize the long-run
growth rate of your individual wealth trajectory, bet this much.
The advice differs from expected-value maximization. The advice is
empirically validated --- investors who follow Kelly outperform
investors who follow expected-value strategies in the long run.

\subsection{The historical neglect}

Despite Kelly's 1956 publication and its substantial subsequent
development by Edward Thorp, William Poundstone, Mark Spitznagel, and
others, the Kelly criterion has not become standard in mainstream
economics or finance. Most textbooks teach expected utility
maximization with a concave utility function (the Bernoulli
framework) and treat Kelly as a special case for ``log utility.''
The teaching obscures the conceptual point: the Kelly criterion is
not a special-case utility function; it is the \emph{correct
criterion} for time-averaged growth, and the failure of standard
expected-utility theory to recommend it is a signal that standard
theory is using the wrong average.

The contemporary ergodicity-economics program, led by Ole Peters at
the London Mathematical Laboratory, has reframed Kelly's result and
its implications. The reframing has not yet been fully absorbed by
mainstream economics, but it is gaining recognition.

\section{Absorbing barriers and bankruptcy}
\label{sec:absorbing-barriers}

A particularly important class of non-ergodic processes involves
absorbing barriers. We treat them in this section because they are
the structural feature of debt-based dynamics that connects the
ergodicity framework to the Islamic prohibition of riba.

\subsection{Definition}

\begin{definition}[Absorbing barrier]
\label{def:absorbing}
An \emph{absorbing barrier} in a stochastic process is a state that,
once entered, cannot be left. The process is ``trapped'' at the
absorbing state for all subsequent times.
\end{definition}

The canonical economic absorbing barrier is bankruptcy. A wealth
process that hits zero (or some negative threshold representing
insolvency) cannot recover; the agent exits the game and stays out.
The wealth dynamics ``before bankruptcy'' and ``after bankruptcy''
are not symmetric; bankruptcy is irreversible.

\subsection{Why absorbing barriers destroy ergodicity}

The presence of an absorbing barrier guarantees that the time
average of any wealth process eventually equals the value at the
absorbing barrier (since the process spends an arbitrarily large
fraction of time there once it has been absorbed). The ensemble
average, by contrast, is dominated by the realizations that have
\emph{not} been absorbed, which can be far above the barrier.

The divergence between time and ensemble averages is therefore
extreme in the presence of absorbing barriers. The ensemble average
gives a reading of how the few survivors are doing; the time
average for any individual gives a reading of where they will end
up, which (if absorption is possible) is the absorbing state.

\subsection{Interest with default risk: a non-ergodic process}

Standard interest-bearing lending has the structural feature that
makes it non-ergodic: in case of default, the lender loses the
principal. The borrower, having defaulted, may also face severe
economic consequences (asset seizure, loss of social standing,
imprisonment in earlier eras, persistent disqualification from
credit in modern eras).

Consider a simplified model. The borrower invests at gross return
$1+r$ each period. With probability $p$, the investment succeeds
and the borrower repays principal plus interest. With probability
$q = 1 - p$, the investment fails and the borrower defaults; their
wealth goes to zero (or below).

For the borrower, the wealth process is non-ergodic with an
absorbing barrier at zero. The expected return per period is
\[
    \mathbb{E}[\text{return}] = p(1 + r) + q \cdot 0 = p(1 + r) - 1.
\]
This is positive if $p(1+r) > 1$, that is, if the success
probability is high enough. The ensemble average grows.

But the time average is dominated by the absorbing barrier. The
probability of survival to time $n$ is $p^n \to 0$ as $n \to
\infty$. With probability 1 in the long run, the borrower is
bankrupt. The time average of wealth is zero.

The divergence between ensemble and time averages here is total: the
ensemble average grows; the time average converges to the absorbing
barrier value. An investor who maximizes the ensemble average will,
with probability 1, eventually be ruined.

\subsection{The Quranic statement on riba}

This is the formal content of the Quranic statement that ``Allah
destroys riba'' (\Quran{2:276}). The statement is not merely moral.
It is mathematically descriptive of what happens to the time average
of wealth in interest-based dynamics with default risk: the time
average converges to the absorbing barrier, regardless of how
favorable the ensemble-average return appears.

We develop this in detail in Chapter 13. The point of the present
section is to establish the mathematical framework: absorbing
barriers destroy ergodicity, and the structural feature of riba ---
the irreversibility of default --- is precisely what creates the
absorbing barrier.

\section{Long-standing economic puzzles dissolved}
\label{sec:dissolved-puzzles}

Several puzzles in the standard economic literature dissolve once
the ergodicity framework is adopted. We treat three briefly.

\subsection{The equity premium puzzle}

Mehra and Prescott (1985) noted that historical stock returns have
exceeded bond returns by an amount that, on standard expected-utility
analysis, requires implausibly high risk aversion to
explain.\footnote{Mehra, R. \& Prescott, E. C. (1985), ``The equity
premium: a puzzle,'' \emph{Journal of Monetary Economics} 15:
145--161.} The standard model predicts a much smaller equity
premium than is observed.

Under ergodicity-economics analysis, the puzzle is reduced. The
relevant comparison is not the ensemble-average returns but the
time-averaged growth rates. Multiplicative wealth dynamics with
high variance (stocks) have a substantial gap between
ensemble-average return and time-averaged growth rate. The
``risk premium'' compensates investors not for risk in some
psychological sense but for the time-ensemble divergence: the
ensemble-average return is higher, but the time-averaged growth
rate is lower; investors must be paid in ensemble-average terms to
accept the lower time-averaged growth.

\subsection{The insurance puzzle}

Why do people buy insurance, given that the insurance company makes
positive expected profit (so the customer's expected return on the
insurance contract is negative)? The standard answer is risk
aversion. The ergodicity answer is more direct: insurance pools
risk, converting non-ergodic individual situations into ergodic
collective ones. The pool experiences the ensemble average; the
individuals, without pooling, would experience the (much lower)
time average. Insurance is rational because it converts the
relevant criterion to one's favor.

\subsection{The cost of inequality}

Standard welfare analysis can struggle to articulate why inequality
is costly. If the total wealth is unchanged, the standard analysis
sees no reduction in welfare from a more unequal distribution.

Ergodicity analysis offers a sharper account. Wealth distributions
have non-ergodic dynamics: small wealth shocks can permanently
disqualify an individual from certain economic activities (because
of capital requirements, credit access, etc.). The ensemble average
of wealth may grow, but the time average for low-wealth individuals
can stagnate or decline if their non-ergodic dynamics keep them
trapped near absorbing barriers. The cost of inequality is then a
real cost, not just a normative judgment.

The Islamic economic prohibitions and prescriptions ---
particularly the package of zakat, riba prohibition, inheritance
partition, and waqf treated in Chapter 17 --- can be read as
addressing exactly this non-ergodic structure of unconstrained
wealth dynamics.

\section{Risk pooling as the rational response}
\label{sec:risk-pooling}

The ergodicity framework leads to a structural conclusion that
mainstream finance has been slow to absorb: risk pooling is not just
risk-aversion accommodation. It is the conversion of a non-ergodic
individual situation into an ergodic collective one.

\subsection{The mechanism}

Consider $N$ agents, each facing a non-ergodic random wealth
dynamics. Each agent's individual wealth has time-averaged growth
rate $g_t < g_e$ (where $g_e$ is the ensemble-average growth rate).
The agents pool their wealth: instead of each holding their own
random wealth, they hold equal shares of the combined pool.

The pool's wealth, by the Law of Large Numbers, has time-averaged
growth approaching the ensemble average $g_e$ as $N$ grows large.
Each agent, holding $1/N$ of the pool, has time-averaged growth
close to $g_e$.

The pool transforms the non-ergodic individual situation into an
ergodic collective one. Each agent's expected long-run growth is
strictly higher in the pool than in the individual case. Pooling is
a Pareto improvement; everyone is better off.

\subsection{The Islamic structures}

The Islamic financial structures that the prohibition of riba was
designed to enable --- \arb{mudarabah} (profit-sharing partnership),
\arb{musharakah} (equity partnership), \arb{takaful} (mutual
insurance) --- are precisely the structures that pool risk in this
way. They are not arbitrary moral preferences; they are the rational
response to non-ergodic wealth dynamics. The prohibition of riba
prevents the absorbing-barrier dynamics that destroy individual
ergodicity; the affirmative recommendations create the pooled
structures that restore it.

This is the deepest connection between ergodicity economics and
Islamic economic claims, and it is what Chapter 13 will develop in
full mathematical detail.

\section{Preview: the connection to riba}
\label{sec:riba-preview}

We end the chapter with a preview of how the ergodicity framework
connects to the Islamic prohibition of riba. The full development is
the subject of Chapter 13.

The Quranic statement is direct: ``Allah destroys riba and increases
charities'' (\Quran{2:276}). The classical reading is moral and
theological. The ergodicity reading is mathematical: in
interest-based lending with default risk, the wealth process is
non-ergodic with an absorbing barrier; the time average of the
borrower's wealth converges to the barrier; the long-run experience
is destruction of wealth. The ensemble average can be positive;
the time average is not. The Quranic statement correctly identifies
the time-averaged outcome.

The classical reading and the ergodicity reading are not in tension.
They are the same insight at different levels of formalism. The
moral and theological reading captures the lived experience of
debt-driven economies (the periodic financial crises, the persistent
inequality, the ruination of borrowers, the destabilization of
communities); the ergodicity reading captures the mathematical
mechanism that produces these outcomes. Both are correct.

The advantage of the ergodicity reading is that it makes the claim
formally derivable. The reader who is told ``riba is forbidden
because Allah forbids it'' has a moral injunction without a
mechanism. The reader who is shown the time-ensemble divergence,
the absorbing-barrier dynamics, and the Jensen inequality
underlying them has a mathematical argument that survives independent
scrutiny. The two readings reinforce each other; neither replaces
the other; together they constitute a stronger case than either
alone.

This is what we mean by Level 3 treatment in the sense of \xref{ch:levels-of-rigor}:
the metaphysical claim corresponds to a real mechanism whose
operation is mathematically demonstrable. The treatment in Chapter
13 carries out the demonstration in full.

\begin{summarybox}[Chapter summary]
\textbf{Probability foundations}: random variables, expectation,
variance, the Strong Law of Large Numbers (which says the sample
average across many independent realizations converges to the
expected value).

\textbf{The hidden assumption}: standard decision theory uses the
expected value (ensemble average) as the criterion for individual
decisions. This is appropriate only if the process is ergodic.

\textbf{Ensemble vs.\ time average}: ensemble average is across
realizations; time average is within a single realization. They are
equal for ergodic processes and differ for non-ergodic ones.

\textbf{The canonical example}: a coin flip with $+50\%$/$-40\%$
multiplicative outcomes has ensemble-average return $+5\%$ per
flip and time-averaged growth $-5.3\%$ per flip. The two are
mathematically distinct; both are correct; they answer different
questions.

\textbf{Theorem on time-ensemble divergence}: for any
multiplicative wealth process with random returns, the time-average
growth rate is strictly less than the ensemble-average growth rate
unless the returns are deterministic. This follows from Jensen's
inequality applied to the concave function $\log$.

\textbf{The Kelly criterion}: $f^* = p - q/b$. Maximizes long-run
time-averaged growth. Differs from expected-value maximization,
which is ruinous for non-ergodic processes.

\textbf{Absorbing barriers}: bankruptcy is the canonical economic
absorbing barrier. Their presence destroys ergodicity completely:
the time average converges to the barrier value, regardless of how
favorable the ensemble-average appears.

\textbf{The connection to riba}: interest with default risk is
non-ergodic with an absorbing barrier. The time-averaged outcome is
destruction of wealth. The Quranic statement that Allah destroys
riba is mathematically descriptive.

\textbf{Risk pooling}: the rational response to non-ergodic
dynamics. \arb{Mudarabah}, \arb{musharakah}, and \arb{takaful} are
precisely the pool structures that convert non-ergodic individual
situations into ergodic collective ones. Pareto improvements.

\textbf{What comes next}: Chapter 7 develops analytical idealism,
the third major resource of Part II.
\end{summarybox}

\cleardoublepage
